\paragraph{} Organisations receive a large volume of semi-structured and unstructured documents such as invoices, contracts, reports, scanned forms, and supporting records. Manual processing of such documents is slow, error-prone, difficult to audit, and hard to scale when users need both structured field extraction and natural-language access to document contents. This project presents an AI-powered document intelligence system that automates document ingestion, text extraction, document classification, key-field extraction, validation, search, and question answering.

\paragraph{} The implemented prototype, named Document Intelligence, uses a React and Vite frontend with a FastAPI backend. Uploaded documents are stored with metadata in SQLite and processed through a modular pipeline. The pipeline detects whether a PDF contains embedded text or requires scanned-document processing. Digital PDFs use direct text extraction, while scanned pages use OCR-based extraction through Tesseract and EasyOCR. A LayoutLMv3-based vision-language component is used for document-type classification and layout-aware feature extraction. A confidence-weighted fusion step combines OCR and layout signals, after which rule-based and LLM-assisted entity extraction identify fields such as dates, amounts, invoice numbers, purchase order numbers, parties, and organisation names. Business validation rules flag suspicious or low-confidence fields for review. The system also indexes document text for keyword search and semantic retrieval.

\paragraph{} A second contribution of the project is an Agentic Retrieval-Augmented Generation workflow for document question answering. Instead of directly passing a question to a language model, the workflow plans the query intent, performs structured-field lookup when possible, retrieves candidate evidence, grades evidence, builds a grounded context, and generates an answer with evidence snippets and confidence. This makes answers more auditable and reduces unsupported responses in comparison with a direct black-box question-answering interface.

\paragraph{} The prototype was evaluated using the benchmark and telemetry hooks implemented in the application. At the current prototype stage, measured metrics include a field extraction F1 score of 0.8, Search Precision@5 of 0.25, Search Mean Reciprocal Rank of 1.0, an estimated throughput of 29.77 documents per hour, and a review rate of 0.152. These results show that the system is suitable as a working proof of concept for intelligent document processing, while also identifying areas for future improvement such as larger benchmark datasets, model fine-tuning, stronger table extraction, production authentication, and cloud-scale deployment.
